\section{Expressiveness of LDNA}
\label{sec:expr}

This section characterises the representational power of the Learnable Dynamic
Neighborhood Aggregation model (LDNA). The reference point is the
one-dimensional Weisfeiler--Leman test (1-WL) of
\citet{weisfeiler1968reduction}, the standard yardstick for message-passing
graph neural networks \citep{xu2019gin,morris2019wl}. All statements are about
the model \emph{class}:
they concern the existence of parameters realising a given behaviour, not about
any particular trained instance. Throughout, $\{\!\{\cdot\}\!\}$ denotes a
multiset (a set with repetitions), and for a finite multiset $M$ over a set $A$
we write $M(a)$ for the multiplicity of $a\in A$.

\subsection{Setup}
\label{sec:expr-setup}

\paragraph{Attributed graphs.}
An attributed graph is $G=(V,E,x,e)$ with node set $V$, directed edge set
$E\subseteq V\times V$ (an undirected edge is the pair of its two orientations),
a node-feature map $x:V\to\Sigma$ and, when present, an edge-feature map
$e:E\to\Xi$, $e_{uv}:=e(u,v)$. We write $N(v)=\{u:(u,v)\in E\}$ for the
in-neighbourhood and $n=|V|$. We make the following standing assumptions.

\begin{itemize}
  \item \textbf{Finite alphabets, bounded size.} The node alphabet
        $\Sigma\subset\mathbb{R}^{d_0}$ and the edge alphabet $\Xi$ are finite,
        and graph sizes are bounded, $n\le N$. This is what makes exact,
        finite-domain constructions available: an MLP with $\mathrm{ReLU}$
        activations can interpolate \emph{any} prescribed function on a finite
        input set exactly (\cref{lem:mlp-interp}), so no appeal to universal
        approximation is needed anywhere below.
  \item \textbf{Minimum degree one.} Every node has at least one neighbour,
        $N(v)\ne\emptyset$. Unlike the previous item this is not universally
        satisfied by the datasets --- isolated nodes do occur in practice (e.g.\
        single-atom fragments in molecular graphs) --- so we invoke it only where
        it is stated. It is used solely by the lower-bound constructions
        (\cref{thm:ldna-lower}, the $\succeq$ direction of
        \cref{prop:delta-charac}, \cref{thm:feature-wl}(a) and
        \cref{thm:canon-strict}), to let a node's own colour be recovered from
        its non-empty aggregation; the upper-bound statements
        (\cref{thm:feature-wl}(b), \cref{prop:invariant-obstruction} and the
        $\preceq$ direction of \cref{prop:delta-charac}) hold for all graphs
        without it. Its necessity is examined in \cref{remark:isolated}.
\end{itemize}

\paragraph{The LDNA layer.}
Nodes carry hidden states $h_v^{(t)}\in\mathbb{R}^{d_t}$, and each node has a
scalar \emph{canonical rank} $r(v)\in[0,1)$ fixed once per graph (defined
below). One LDNA layer (the module \texttt{LDNAConv}) computes, for every
directed edge $u\to v$ with $u\in N(v)$, a pre-message, a per-channel gate of
the rank gap, and a post-transform of the gated sum:
\begin{align}
  m_{uv} &= \operatorname{pre}\!\big([\,h_v^{(t)},\,h_u^{(t)},\,e_{uv}\,]\big),
  \label{eqn:expr-pre}\\
  \delta_{uv} &= r(u)-r(v),\qquad
  g(\delta) = \mathbf{1} + \operatorname{MLP}_g\!\big([\,\delta,\ |\delta|,\ \operatorname{sign}\delta\,]\big),
  \label{eqn:expr-gate}\\
  h_v^{(t+1)} &= \operatorname{post}\!\Big(\ \sum_{u\in N(v)} g(\delta_{uv})\odot m_{uv}\ \Big),
  \label{eqn:expr-update}
\end{align}
where $\odot$ is the Hadamard product, $\operatorname{pre}$ and
$\operatorname{post}$ are MLPs (the term $e_{uv}$ is omitted when there are no
edge features), and $\operatorname{MLP}_g$ is a two-layer network
$\mathrm{Linear}\to\mathrm{ReLU}\to\mathrm{Linear}$ whose output has the channel
dimension of $m_{uv}$. The output layer of $\operatorname{MLP}_g$ is initialised
to zero, so $g\equiv\mathbf{1}$ is exactly realisable (and is the value at
initialisation): with $g\equiv\mathbf{1}$ the layer is a plain summing
message-passing layer.

\paragraph{The LDNA network.}
The wrapper (\texttt{LDNA}) applies a node encoder and, if present, an edge
encoder; then a stack of $L$ blocks, each a layer
\labelcref{eqn:expr-pre,eqn:expr-gate,eqn:expr-update} followed by BatchNorm, a
pointwise activation and dropout, with an optional per-layer residual
$h_v^{(t+1)}\mathrel{+}= W h_v^{(t)}$ ($W$ the identity where the layer preserves
width, a learnable linear map where it changes width); then a permutation-%
invariant readout in
$\{\textsf{add},\textsf{mean},\textsf{max},\textsf{attention}\}$ (\textsf{add}
$=$ sum, \textsf{mean} $=$ average, \textsf{max} $=$ coordinatewise maximum,
\textsf{attention} $=$ a softmax-weighted sum, PyG's
\texttt{AttentionalAggregation}) or the order-sensitive \textsf{gru} readout (a
GRU over the sorted node sequence); and finally a prediction MLP. We treat the
\textsf{gru} readout separately in \cref{remark:gru}; the four listed readouts
are all functions of the multiset $\{\!\{h_v^{(L)}:v\in V\}\!\}$.

\begin{remark}[Inference-time normalisation]
\label{remark:bn}
We analyse the network at inference. Dropout is then inactive, and BatchNorm is
a fixed affine map $z\mapsto\gamma\odot(z-\mu)/\sigma+\beta$, which can be
absorbed into the adjacent MLP; wherever a construction needs it we take this
affine map to be the identity. We do not mention BatchNorm or dropout again.
\end{remark}

\paragraph{Canonical rank.}
The rank $r$ is supplied to \cref{eqn:expr-gate} in one of two modes. In both,
only within-graph gaps $\delta_{uv}=r(u)-r(v)$ are ever used (all edges are
intra-graph), so any per-graph additive offset in $r$ is immaterial.

\emph{Feature mode.} Nodes are first put into a fixed data-side canonical order
by sorting each graph lexicographically by its raw feature rows
(\texttt{sort\_graph}); this fixes a total order $<_{\mathrm{lex}}$ on the finite
alphabet $\Sigma$. Let $X_G=\{x_v:v\in V\}\subseteq\Sigma$ be the \emph{set} of
distinct rows occurring in $G$ and
\begin{equation}
  \rho_G(x)\;=\;\big|\{\,x'\in X_G:\ x'<_{\mathrm{lex}}x\,\}\big|
  \label{eqn:expr-dense-rank}
\end{equation}
the tie-aware \emph{dense} rank of a row (nodes sharing a row share a rank).
The feature-mode rank is
\begin{equation}
  r(v)\;=\;\rho_G(x_v)\,/\,n .
  \label{eqn:expr-feature-rank}
\end{equation}
Normalisation is by $n=|V|$, the number of nodes, not the number of distinct
rows.

\emph{Position mode.} When the data already supplies a canonical total order of
the nodes by construction (for instance the abstract-syntax-tree depth-first
order of \texttt{ogbg-code2}), $r(v)=p(v)/n$, where $p(v)\in\{0,\dots,n-1\}$ is
the position of $v$ in that order.

\paragraph{The reference test.}
The 1-WL colour refinement with edge colours assigns
\begin{equation}
  c_v^{(0)}=x_v,\qquad
  c_v^{(t+1)}=\operatorname{HASH}\!\Big(c_v^{(t)},\ \{\!\{\,(c_u^{(t)},e_{uv}):u\in N(v)\,\}\!\}\Big),
  \label{eqn:expr-wl}
\end{equation}
with $\operatorname{HASH}$ injective; colours range over a finite set at each
round. The round-$(t{+}1)$ partition refines the round-$t$ partition, and it is
monotone: two nodes with distinct round-$t$ colours have distinct round-$(t{+}1)$
colours. We say $G$ and $H$ are \emph{1-WL-indistinguishable}, $G\sim_{1\text{-WL}}H$,
if for every $t$ their colour histograms agree,
$\{\!\{c_v^{(t)}:v\in V_G\}\!\}=\{\!\{c_w^{(t)}:w\in V_H\}\!\}$; this forces
$|V_G|=|V_H|$. Because refinement is monotone, if the histograms differ at some
round $t^\ast$ they differ at every round $t\ge t^\ast$; in particular they
differ at the stable round $T\le N-1$ at which refinement terminates on all
graphs of size $\le N$.

\begin{definition}[Distinguishing power]
\label{def:expr}
For a model class $\mathcal{M}$ (all parameters and widths), $\mathcal{M}$
\emph{distinguishes} a pair $(G,H)$ if there exist parameters $\theta$ with
$\mathcal{M}_\theta(G)\ne\mathcal{M}_\theta(H)$. We say $\mathcal{M}$ is
\emph{at least as expressive as} 1-WL, $\mathcal{M}\succeq 1\text{-WL}$, if it
distinguishes every 1-WL-distinguishable pair; $\mathcal{M}$ is \emph{strictly
more expressive} than 1-WL if in addition it distinguishes some pair with
$G\sim_{1\text{-WL}}H$.
\end{definition}

\paragraph{Two constructive lemmas.}
The proofs use only the following two elementary facts; neither uses universal
approximation.

\begin{lemma}[Exact finite interpolation]
\label{lem:mlp-interp}
Let $S\subset\mathbb{R}^d$ be finite and $\psi:S\to\mathbb{R}^m$ arbitrary.
There is a two-layer $\mathrm{ReLU}$ network
$F(z)=W_2\,\sigma(W_1z+b_1)+b_2$ ($\sigma=\mathrm{ReLU}$) with $F|_S=\psi$.
\end{lemma}

\begin{proof}
Choose $w\in\mathbb{R}^d$ with the scalars $t_s:=\langle w,s\rangle$ pairwise
distinct over $s\in S$; the excluded $w$ form the finite union of hyperplanes
$\{w:\langle w,s-s'\rangle=0\}$, $s\ne s'$, so such $w$ exist. Order
$S=\{s_1,\dots,s_K\}$ so that $t_1<\dots<t_K$, and write $y_k=\psi(s_k)$. It
suffices to build, for each output coordinate $j$, a continuous piecewise-linear
$f_j:\mathbb{R}\to\mathbb{R}$ with $f_j(t_k)=(y_k)_j$ and to set
$F(z)_j=f_j(\langle w,z\rangle)$. Writing
$a_k=\big((y_{k+1})_j-(y_k)_j\big)/(t_{k+1}-t_k)$ for the slope on
$[t_k,t_{k+1}]$ ($k=1,\dots,K-1$) and setting $a_0:=0$, define
\[
  f_j(t)\;=\;(y_1)_j\;+\;\sum_{k=1}^{K-1} (a_k-a_{k-1})\,\sigma(t-t_k) .
\]
For $t\le t_1$ every $\sigma(t-t_k)$ vanishes, so $f_j(t)=(y_1)_j$; in particular
$f_j(t_1)=(y_1)_j$. On the interval $(t_k,t_{k+1})$ exactly the units
$\sigma(t-t_1),\dots,\sigma(t-t_k)$ are active, so the slope of $f_j$ telescopes
to $\sum_{i=1}^{k}(a_i-a_{i-1})=a_k$; hence there $f_j$ is affine with slope
$a_k$, and $f_j(t_{k+1})=f_j(t_k)+a_k(t_{k+1}-t_k)=(y_{k+1})_j$, so by induction
on $k$ the map $f_j$ interpolates the points $(t_k,(y_k)_j)$. Each of the $K-1$
hidden units computes $\sigma(\langle w,z\rangle-t_k)$ (a single $\mathrm{ReLU}$
of an affine form), and the outer layer takes the stated linear combination with
the constant $(y_1)_j$ in the output bias, so $F$ is a two-layer $\mathrm{ReLU}$
network.
\end{proof}

\begin{lemma}[Sum-injective encodings]
\label{lem:sum-inj}
Let $A$ be finite, $|A|=K$, and fix a bijection $\iota:A\to\{1,\dots,K\}$. The
one-hot map $\chi(a)=e_{\iota(a)}\in\mathbb{R}^K$ satisfies, for all finite
multisets $M,M'$ over $A$,
$\sum_{a\in M}\chi(a)=\sum_{a\in M'}\chi(a)\iff M=M'$. Moreover, for multisets of
size $\le N$ the scalar $\sum_{a\in M}(N+1)^{\iota(a)-1}$ is injective in $M$.
\end{lemma}

\begin{proof}
$\sum_{a\in M}\chi(a)$ is the count vector $\big(M(a)\big)_{a\in A}$, which
determines $M$ and conversely. For the scalar version, the multiplicities
$M(a)\le N<N+1$ are the digits of a base-$(N+1)$ integer, whose representation is
unique.
\end{proof}

The $\operatorname{pre}$ and $\operatorname{post}$ MLPs below are taken with at
least one hidden layer (i.e.\ \texttt{num\_pre\_layers},
\texttt{num\_post\_layers}$\ \ge 2$), a valid configuration; \cref{lem:mlp-interp}
is then applied to realise them exactly on their finite domains. The node and
edge encoders are the implemented ones and are \emph{not} assumed to be such
MLPs: they need only be \emph{injective} on the finite alphabet, which every
implemented family provides --- a single embedding table (choose distinct rows),
an additive multi-column encoder (\texttt{AtomEncoder}, a sum of per-column
embeddings placed on disjoint coordinate blocks, so the sum is a concatenation),
or a scalar $\mathrm{Linear}(1,d)$ (any nonzero weight). Finally, in the
implementation the $\operatorname{pre}$ and gate output width is tied to the
layer's input width, so ``message width'' below refers to the layer's hidden
width, which is free (all widths are quantified over in \cref{def:expr}).

\subsection{A lower bound in every rank mode}
\label{sec:expr-lower}

\begin{proposition}[Lower bound]
\label{thm:ldna-lower}
Under minimum degree one, with the \textsf{add} readout, LDNA is at least as
expressive as 1-WL, $\text{LDNA}\succeq 1\text{-WL}$, in every rank mode.
\end{proposition}

\begin{proof}
Set the gate to the constant $g\equiv\mathbf{1}$ by zeroing the output layer of
$\operatorname{MLP}_g$ (\cref{eqn:expr-gate}); the rank then plays no role and the
argument is identical in every mode. The layer reduces to
$h_v^{(t+1)}=\operatorname{post}\big(\sum_{u\in N(v)}\operatorname{pre}([h_v^{(t)},h_u^{(t)},e_{uv}])\big)$.
Let $C_t$ be the finite set of round-$t$ 1-WL colours.

We claim that there are parameters and, for each $t$, a map
$\alpha_t:C_t\to\mathbb{R}^{d_t}$ with $h_v^{(t)}=\alpha_t(c_v^{(t)})$ for every
node of every graph of size $\le N$, where $\alpha_0$ is injective and, for
$t\ge 1$, $\alpha_t$ is one-hot. (One layer already produces a one-hot code, and
the readout argument below runs at some $L\ge 1$.)

\emph{Base.} $c_v^{(0)}=x_v\in\Sigma$, and the node encoder gives
$h_v^{(0)}=\alpha_0(x_v)$ with $\alpha_0$ injective on the finite $\Sigma$ (any of
the implemented encoders, as above). Injectivity of $\alpha_0$ is all that the
step below uses at $t=0$.

\emph{Step.} Assume $h_v^{(t)}=\alpha_t(c_v^{(t)})$ with $\alpha_t$ injective
(and one-hot when $t\ge 1$). Since $\alpha_t$ is invertible on its image, from
$[h_v^{(t)},h_u^{(t)},e_{uv}]$ one can read off the triple
$(c_v^{(t)},c_u^{(t)},e_{uv})$, an element of $C_t\times C_t\times\Xi$. Choose
$\operatorname{pre}$, via \cref{lem:mlp-interp}, to output the one-hot
$\chi(c_v^{(t)},c_u^{(t)},e_{uv})$ of this triple. Because $c_v^{(t)}$ is
constant over $u\in N(v)$, the aggregated vector
\[
  J_v \;=\; \sum_{u\in N(v)} \chi\big(c_v^{(t)},c_u^{(t)},e_{uv}\big)
\]
is supported on the block indexed by $c_v^{(t)}$, and within that block it is,
by \cref{lem:sum-inj}, the histogram of
$\{\!\{(c_u^{(t)},e_{uv}):u\in N(v)\}\!\}$. Thus $J_v$ encodes the pair
$\big(c_v^{(t)},\{\!\{(c_u^{(t)},e_{uv})\}\!\}\big)$ injectively; the aggregation
is non-empty because no node is isolated. By \cref{eqn:expr-wl} and injectivity of
$\operatorname{HASH}$, this pair is in bijection with $c_v^{(t+1)}$, so $J_v$
takes a distinct value for each colour $c_v^{(t+1)}\in C_{t+1}$, over a finite
set of possible values. Choose $\operatorname{post}$, via \cref{lem:mlp-interp},
to send each such $J_v$ to the one-hot $\alpha_{t+1}(c_v^{(t+1)})$. One-hot
codes are non-negative, so the intervening $\mathrm{ReLU}$ acts as the identity
(\cref{remark:bn}). Take the widths strictly increasing so that every residual
projection is a learnable map, set to zero; the residual connections are thus
present but inert here. This proves the claim.

\emph{Graph level.} Run $L=\max(T,1)$ layers, where $T\le N-1$ is the stable
round; then $L\ge 1$, so $h_v^{(L)}=\alpha_L(c_v^{(L)})$ is one-hot, and the
\textsf{add} readout returns $\sum_{v}\alpha_L(c_v^{(L)})$, the histogram of
round-$L$ colours, which by \cref{lem:sum-inj} is injective in the multiset
$\{\!\{c_v^{(L)}\}\!\}$. If $G,H$ are 1-WL-distinguishable they differ at some
round, hence (by monotonicity) at round $L$, so the two readouts differ; an
identity (or any injective) prediction head preserves the difference. Hence LDNA
distinguishes $(G,H)$.
\end{proof}

\begin{remark}[Isolated nodes and the necessity of minimum degree one]
\label{remark:isolated}
The hypothesis $N(v)\ne\emptyset$ is essential to \cref{thm:ldna-lower}, not
cosmetic, and isolated nodes do occur in practice (for instance single-atom
fragments in molecular graphs). One might hope the residual could rescue them,
but it cannot, because of the layer ordering in the implementation: the
activation is applied to the convolution output \emph{before} the residual is
added, and for an isolated node the aggregation is $J_v\equiv 0$ for every choice
of parameters. Its update therefore collapses to an affine recurrence
$h_v^{(t+1)}=b_t+W_t h_v^{(t)}$, with $b_t=\mathrm{ReLU}(\mathrm{BN}(\operatorname{post}(0)))$
a constant and $W_t$ the residual projection, so $h_v^{(L)}$ is an \emph{affine}
function of $\text{enc}(x_v)$. Take moreover the node encoder to be affine in a
scalar feature --- a $\mathrm{Linear}(1,d)$ encoder, a supported configuration ---
so $h_v^{(L)}=c+Ax_v$ for a fixed vector $c$ and matrix $A$. Then on any two
graphs consisting only of isolated nodes the \textsf{add} readout returns
$\sum_v(c+Ax_v)=|V|\,c+A\sum_v x_v$, which depends only on the node count and the
\emph{sum} of features: the feature multisets $\{\!\{1,3\}\!\}$ and
$\{\!\{2,2\}\!\}$ collide, though 1-WL separates them. Thus for the class with an
affine-in-a-scalar encoder and the \textsf{add} readout, minimum degree one
cannot be dropped from \cref{thm:ldna-lower}.
\end{remark}

\begin{remark}[Cardinality-blind readouts]
\label{remark:readout-cardinality}
The proposition uses \textsf{add}, which the class includes. The \textsf{mean}
and \textsf{attention} readouts are normalised (softmax weights sum to one), so
they see only the \emph{distribution} of final embeddings, not their multiset
with multiplicities. Consequently they cannot separate a graph $G$ from the
disjoint union $G\sqcup G$: the colour multiset of $G\sqcup G$ is that of $G$
with every multiplicity doubled, so \textsf{mean} and \textsf{attention} return
identical values on $G$ and $G\sqcup G$, whereas 1-WL distinguishes them (their
colour histograms, and indeed their sizes, differ). This is why the lower bound
is stated for \textsf{add}; cf.\ the analogous discussion of the mean aggregator
by \citet{xu2019gin}.
\end{remark}

\subsection{Exact characterisation via a rank-augmented test}
\label{sec:expr-char}

Fix a rank function $r$. Define the \emph{rank-augmented graph} $G^r$ to have the
same nodes, edges and features as $G$, with each directed edge $u\to v$ carrying
the augmented colour $(e_{uv},\delta_{uv})$, $\delta_{uv}=r(u)-r(v)$. Let
$1\text{-WL}^r$ denote edge-coloured 1-WL run on $G^r$: it is
\cref{eqn:expr-wl} with the neighbour term replaced by
$\{\!\{(c_u^{(t)},e_{uv},\delta_{uv}):u\in N(v)\}\!\}$.

\begin{proposition}[Characterisation]
\label{prop:delta-charac}
For any fixed rank $r$, LDNA with rank $r$ has exactly the distinguishing power
of $1\text{-WL}^r$; the $\succeq$ direction assumes minimum degree one.
\end{proposition}

\begin{proof}
\emph{Upper bound ($\preceq$).} Each LDNA message
$g(\delta_{uv})\odot m_{uv}$ depends only on
$(h_v^{(t)},h_u^{(t)},e_{uv},\delta_{uv})$, so a layer is a message-passing
update on the edge-coloured graph $G^r$. By induction on $t$ there is a map
$\Phi_t$ (parameter-dependent, but the same function for all graphs) with
$h_v^{(t)}=\Phi_t\big(c_v^{(t),r}\big)$, where $c^{(\cdot),r}$ are the
$1\text{-WL}^r$ colours: the base is the encoder, and the step holds because
\cref{eqn:expr-pre,eqn:expr-gate,eqn:expr-update} turn
$\big(c_v^{(t),r},\{\!\{(c_u^{(t),r},e_{uv},\delta_{uv})\}\!\}\big)$ into
$h_v^{(t+1)}$, and that pair is determined by $c_v^{(t+1),r}$. Every
permutation-invariant readout is a function of the final colour histogram, so
graphs with equal $1\text{-WL}^r$ histograms at all rounds receive equal
outputs.

\emph{Lower bound ($\succeq$).} Assuming minimum degree one, we adapt
\cref{thm:ldna-lower} to also encode $\delta$. The gate is per-channel and
depends on $\delta$, while
$\operatorname{pre}$ depends on $(h_v,h_u,e)$ but not on $\delta$; their Hadamard
product lets us form a \emph{joint} one-hot of the quadruple
$(c_v^{(t)},c_u^{(t)},e_{uv},\delta_{uv})$. Concretely, $\delta$ ranges over a
finite set $D$ (there are finitely many graphs of size $\le N$, hence finitely
many gap values; for the model's own rank modes $D\subseteq\{k/n:|k|<n\le N\}$);
index the message channels by
$(a,b,\varepsilon,j)$ with $(a,b,\varepsilon)\in C_t\times C_t\times\Xi$ and
$j$ indexing $D$. Let $\operatorname{pre}$ output, on channel
$(a,b,\varepsilon,j)$, the indicator
$\mathbf{1}[(c_v^{(t)},c_u^{(t)},e_{uv})=(a,b,\varepsilon)]$ (independent of $j$),
and let the gate output, on that channel,
$\mathbf{1}[\delta_{uv}=\delta_j]$ (independent of $(a,b,\varepsilon)$): the map
$\delta\mapsto[\delta,|\delta|,\operatorname{sign}\delta]$ is injective on the
finite $D$, so \cref{lem:mlp-interp} realises this indicator (minus one) with
$\operatorname{MLP}_g$. The product on channel $(a,b,\varepsilon,j)$ is
$\mathbf{1}[(c_v^{(t)},c_u^{(t)},e_{uv},\delta_{uv})=(a,b,\varepsilon,\delta_j)]$,
a one-hot of the quadruple. Summing over $u\in N(v)$ (with $c_v^{(t)}$ constant)
gives, exactly as in \cref{thm:ldna-lower}, an injective encoding of
$\big(c_v^{(t)},\{\!\{(c_u^{(t)},e_{uv},\delta_{uv})\}\!\}\big)$, which by
injectivity of $\operatorname{HASH}$ is in bijection with $c_v^{(t+1),r}$;
$\operatorname{post}$ maps it to the one-hot $\alpha_{t+1}(c_v^{(t+1),r})$. Thus
LDNA simulates $1\text{-WL}^r$, and the \textsf{add} readout recovers its colour
histogram, as before.
\end{proof}

\subsection{Feature mode collapses to 1-WL}
\label{sec:expr-feature}

\begin{lemma}[Rank rigidity in feature mode]
\label{lem:rank-rigid}
If $G\sim_{1\text{-WL}}H$ then
\begin{enumerate}
  \item[(i)] $n_G=n_H=:n$;
  \item[(ii)] the multisets of raw feature rows of $G$ and $H$ coincide, hence
        the sets $X_G=X_H$ of distinct rows coincide and inherit the same order
        $<_{\mathrm{lex}}$;
  \item[(iii)] $\rho_G=\rho_H=:\rho$ as functions, and the feature-mode rank is
        $r(v)=\rho(x_v)/n$ in both graphs --- one fixed function of a node's own
        raw feature row, shared by $G$ and $H$.
\end{enumerate}
\end{lemma}

\begin{proof}
$G\sim_{1\text{-WL}}H$ requires equal round-$0$ histograms,
$\{\!\{x_v:v\in V_G\}\!\}=\{\!\{x_w:w\in V_H\}\!\}$. Equal multiset cardinalities
give (i); equal multisets have equal underlying sets $X_G=X_H$, and
$<_{\mathrm{lex}}$ is a fixed total order on $\Sigma$ that restricts identically
to the common set, giving (ii). By \cref{eqn:expr-dense-rank}, $\rho_G$ depends only on
$(X_G,<_{\mathrm{lex}})$, so $\rho_G=\rho_H=\rho$; with the common $n$ from (i),
\cref{eqn:expr-feature-rank} gives (iii).
\end{proof}

\begin{theorem}[Feature mode equals 1-WL]
\label{thm:feature-wl}
In feature mode:
\begin{enumerate}
  \item[(a)] under minimum degree one, with the \textsf{add} readout,
        $\text{LDNA}\succeq 1\text{-WL}$;
  \item[(b)] for every parameter setting and every permutation-invariant readout
        ($\textsf{add}$, $\textsf{mean}$, $\textsf{max}$, $\textsf{attention}$),
        $G\sim_{1\text{-WL}}H$ implies identical outputs on $G$ and $H$.
\end{enumerate}
Hence feature-mode LDNA is exactly as expressive as 1-WL.
\end{theorem}

\begin{proof}
Part (a) is \cref{thm:ldna-lower}, which is rank-mode-agnostic. For part (b),
fix $G\sim_{1\text{-WL}}H$ and any parameters. By \cref{lem:rank-rigid} the gap
\begin{equation}
  \delta_{uv}=\frac{\rho(x_u)-\rho(x_v)}{n}
  \label{eqn:expr-delta-folded}
\end{equation}
is one fixed function of the raw endpoint rows, common to $G$ and $H$.

We show by induction on $t$ that there is a function $\phi_t$ (parameter-%
dependent, common to both graphs) with $h_v^{(t)}=\phi_t(c_v^{(t)})$ for all
nodes of $G$ and $H$. The base is $h_v^{(0)}=\text{encoder}(x_v)$ with
$c_v^{(0)}=x_v$, so $\phi_0=\text{encoder}$. For the step, recall that
$c_v^{(t)}$ refines $c_v^{(0)}=x_v$, so $c_v^{(t)}$ determines $x_v$ and hence
$\rho(x_v)$; likewise for each neighbour. Therefore every ingredient of
\cref{eqn:expr-pre,eqn:expr-gate,eqn:expr-update} --- the messages
$m_{uv}=\operatorname{pre}([\phi_t(c_v^{(t)}),\phi_t(c_u^{(t)}),e_{uv}])$, the
gate $g(\delta_{uv})$ with $\delta_{uv}$ read off from $(c_v^{(t)},c_u^{(t)})$
via \cref{eqn:expr-delta-folded}, the sum over $u\in N(v)$, and $\operatorname{post}$
--- is determined by
$\big(c_v^{(t)},\{\!\{(c_u^{(t)},e_{uv}):u\in N(v)\}\!\}\big)$, which by
injectivity of $\operatorname{HASH}$ is determined by $c_v^{(t+1)}$. The residual
$W\phi_t(c_v^{(t)})$, the affine BatchNorm and the pointwise activation act
node-wise and preserve this dependence (each is a function of $c_v^{(t+1)}$,
which determines $c_v^{(t)}$). Set $\phi_{t+1}$ accordingly.

At the graph level, $G\sim_{1\text{-WL}}H$ gives equal final colour histograms,
hence equal multisets $\{\!\{h_v^{(T)}\}\!\}=\{\!\{\phi_T(c_v^{(T)})\}\!\}$. The
readouts $\textsf{add}$, $\textsf{mean}$ (the sizes agree by
\cref{lem:rank-rigid}(i)), $\textsf{max}$ and $\textsf{attention}$ are all
functions of this multiset, so the graph embeddings agree; the prediction MLP is
pointwise, so the outputs agree.

Equivalently, via \cref{prop:delta-charac}: by \cref{eqn:expr-delta-folded} the
augmented edge colour $(e_{uv},\delta_{uv})$ is a function of the endpoints'
round-$0$ colours, so $1\text{-WL}^r$ and $1\text{-WL}$ induce the same
partition, and the characterisation collapses to plain 1-WL. Combining (a) and
(b) proves the claim.
\end{proof}

\begin{corollary}
\label{cor:arg-config}
Consider attributed regular graphs --- graphs whose nodes split into classes
$V_1,\dots,V_k$, all nodes of class $V_i$ sharing the attribute $C_i$, such that
every node of class $V_i$ has exactly $r_{ij}$ neighbours in class $V_j$. Encode
this as the configuration $\{(C_i,C_j,r_{ij}):1\le i,j\le k\}$, which records the
attributes as well as the counts. Then any two such graphs without edge features
(or with edge features constant on each ordered pair of classes) that have equal
configurations and equal class sizes are 1-WL-indistinguishable, and hence
receive identical feature-mode LDNA outputs for every parameter setting.
\end{corollary}

\begin{proof}
The partition into classes is equitable and each class carries a fixed attribute:
the counts $r_{ij}$ and the attributes $C_i$ are constant across each class, and
the edge colours (absent, or constant on each ordered class pair) add nothing.
Hence 1-WL initialised at the attributes stabilises at this partition, and each
round-$t$ colour is a function of a node's class alone, determined by the
configuration. Two such graphs with equal configurations and equal class sizes
therefore have equal round-$0$ histograms and equal colour histograms at every
round, i.e.\ $G\sim_{1\text{-WL}}H$. Apply \cref{thm:feature-wl}(b).
\end{proof}

\subsection{Invariant tie-breaking does not help robustly}
\label{sec:expr-invariant}

\Cref{thm:feature-wl} shows that a rank read off the raw features is folded away.
One might hope to strengthen the rank by sorting on a richer node quantity. The
next proposition shows that any \emph{isomorphism-invariant} choice of rank
remains bounded by 1-WL, and fails outright on the classical hard instances.

\begin{proposition}[Obstruction for invariant ranks]
\label{prop:invariant-obstruction}
\begin{enumerate}
  \item[(i)] Let $r$ be obtained by sorting nodes on any node invariant $\iota$
        that is a function of the stable 1-WL colour $\kappa$ (for instance the
        degree, or $\kappa$ itself). Then LDNA with rank $r$ is bounded by 1-WL.
  \item[(ii)] Let $r$ be obtained by sorting on any isomorphism-invariant node
        function (possibly stronger than 1-WL, e.g.\ per-node subgraph counts).
        Then for any pair of vertex-transitive graphs with
        $G\sim_{1\text{-WL}}H$, LDNA with rank $r$ cannot distinguish $G$ from
        $H$.
\end{enumerate}
\end{proposition}

\begin{proof}
\emph{(i).} Since $\iota$ is a function of $\kappa$, and the dense rank
\cref{eqn:expr-dense-rank} of $\iota$-values divided by the common size is, for
$G\sim_{1\text{-WL}}H$, a fixed function of $\kappa$ (the multiset of $\kappa$,
hence of $\iota$, agrees across $G$ and $H$), the gap
$\delta_{uv}=f(\kappa(u),\kappa(v))$ is determined by the endpoints' stable
colours. By \cref{prop:delta-charac}, LDNA with $r$ is bounded by
$1\text{-WL}^r$. We claim $1\text{-WL}^r$ and $1\text{-WL}$ induce the same
partition. On one hand $1\text{-WL}^r$ refines $1\text{-WL}$. On the other, the
stable partition $\kappa$ is equitable for $G$, and the added edge colours
$\delta=f(\kappa(u),\kappa(v))$ are constant on each ordered pair of classes;
hence for a node in class $C$ the multiset
$\{\!\{(\kappa(u),e_{uv},\delta_{uv}):u\in N(v)\}\!\}$ is a fixed function of the
un-augmented multiset $\{\!\{(\kappa(u),e_{uv})\}\!\}$, which is constant across
$C$. Thus $\kappa$ is equitable for $G^r$ as well, so $1\text{-WL}^r$ does not
refine $\kappa$ and equals it. To pass from equal partitions to equal
graph-level histograms, check by induction on $t$ that every $1\text{-WL}^r$
colour $c_v^{(t),r}$ is the image of $\kappa(v)$ under a map shared by $G$ and
$H$. Base: $c_v^{(0),r}=x_v$, and $\kappa$ refines the raw feature, so $x_v$ is a
function of $\kappa(v)$. Step:
\[
  c_v^{(t+1),r}=\operatorname{HASH}\!\Big(c_v^{(t),r},\ \{\!\{(c_u^{(t),r},e_{uv},\delta_{uv}):u\in N(v)\}\!\}\Big);
\]
by the induction hypothesis and $\delta_{uv}=f(\kappa(u),\kappa(v))$ with $f$
shared (the dense-rank map is shared, as above), the augmented neighbour multiset
is a shared-map image of $\mu_v:=\{\!\{(\kappa(u),e_{uv}):u\in N(v)\}\!\}$, and
$\mu_v$ is itself a fixed function of $\kappa(v)$ because $\kappa$ is equitable;
hence $c_v^{(t+1),r}$ is a shared function of $\kappa(v)$. Consequently each
per-round $1\text{-WL}^r$ colour histogram is the image, under a shared map, of
the corresponding $\kappa$-histogram. The $\kappa$-histograms agree across $G$
and $H$ (since $G\sim_{1\text{-WL}}H$ fixes every round histogram, including the
stable one), so the $1\text{-WL}^r$ histograms agree too. Therefore LDNA with
$r$ distinguishes no pair that 1-WL cannot.

\emph{(ii).} If $\iota$ is isomorphism-invariant and $G$ is vertex-transitive,
then for any $u,v$ there is an automorphism $\pi$ with $\pi(u)=v$, whence
$\iota(u)=\iota(v)$; so $\iota$ is constant and all nodes are tied. Then
$r$ is constant, $\delta_{uv}\equiv 0$ on every edge, and the gate takes the
single value $g(0)=\mathbf{1}+\operatorname{MLP}_g([0,0,0])$, a constant that
folds into $\operatorname{post}$. LDNA thereby degenerates to the gate-free model
of \cref{thm:ldna-lower}; equivalently $G^r=G$ (all $\delta=0$), so by
\cref{prop:delta-charac} LDNA with $r$ is bounded by 1-WL on this pair and cannot
distinguish $G\sim_{1\text{-WL}}H$. The classical 1-WL-hard pairs are
vertex-transitive --- for example the $4\times 4$ rook's graph versus the
Shrikhande graph, and the circular-skip-link (CSL) family --- so invariant-based
tie-breaking provably fails on them.
\end{proof}

\subsection{Canonical rank and strict separation}
\label{sec:expr-canonical}

The obstruction in \cref{prop:invariant-obstruction} is that an
isomorphism-invariant rank must respect automorphisms and so ties every orbit.
Breaking the symmetry requires a rank that is invariant across graphs but not
across the automorphisms of a single graph --- exactly what a canonical
labelling provides.

\begin{definition}[Injective canonical rank]
\label{def:canonical-rank}
A canonical labelling algorithm assigns to each attributed graph $G$ a bijection
$\pi_G:V\to\{0,\dots,n-1\}$ so that isomorphic graphs receive identical labelled
graphs (a canonical form): for every isomorphism $f:G\to H$, relabelling $G$ by
$\pi_G$ and $H$ by $\pi_H$ yields the same ordered attributed graph. The induced
canonical rank is $r(v)=\pi_G(v)/n$; it is injective on the nodes of $G$.
\end{definition}

\begin{remark}[Properties of the canonical rank]
\label{remark:canonical}
Three points attach to \cref{def:canonical-rank}.
\emph{Well-definedness under automorphisms.} A canonizer may choose any
representative labelling within an automorphism orbit; all such choices give the
same canonical form, so the end-to-end model output is an exact, deterministic
graph invariant.
\emph{Realisation.} The \texttt{position} rank mode implements this exactly when
the data supply a canonical order by construction (the AST depth-first order of
\texttt{ogbg-code2}). Alternatively the feature lexsort is completed into a full
canonization by individualization--refinement \citep{mckay2014nauty}: the lexsort
is the initial colour classification, colour refinement \emph{is} 1-WL, and what
is added is a canonical individualization of the remaining ties.
\emph{Complexity.} Graph canonization has no known polynomial worst-case
algorithm --- the best bound is quasi-polynomial \citep{babai2016quasipoly} ---
but is fast in practice on benchmark graph classes \citep{mckay2014nauty}; it is
a one-time, data-side preprocessing step, in the same pipeline position as the
existing lexsort.
\end{remark}

\begin{theorem}[Strictness under canonical rank]
\label{thm:canon-strict}
With an injective canonical rank, under minimum degree one, and the \textsf{add}
readout, LDNA is strictly more expressive than 1-WL.
\end{theorem}

\begin{proof}
That $\text{LDNA}\succeq 1\text{-WL}$ is \cref{thm:ldna-lower}, which holds in
every rank mode. It remains to exhibit a 1-WL-indistinguishable pair that LDNA
separates.

Let $G_1=C_6$ be the $6$-cycle and $G_2=C_3\sqcup C_3$ two disjoint triangles,
all node features equal to a single symbol and no edge features. Both are
$2$-regular on six nodes (so minimum degree one holds) with identical initial
colours. By induction, at every
1-WL round all six nodes of either graph share one colour: at round $0$ this is
the single symbol, and if all nodes share colour $a$ at round $t$, then each
node's neighbour multiset is $\{\!\{a,a\}\!\}$, so $\operatorname{HASH}$ returns
one common colour at round $t+1$. The round-$t$ histograms are therefore
$\{\!\{a_t^{\,(6)}\}\!\}$ for both graphs and agree at every $t$; hence
$G_1\sim_{1\text{-WL}}G_2$.

A canonizer is free to fix one representative labelling per isomorphism class.
Choose $\pi(C_6)$ to be the cyclic order (positions $0,1,2,3,4,5$ along the
cycle) and $\pi(C_3\sqcup C_3)$ to place the two triangles on
$\{0,1,2\}$ and $\{3,4,5\}$. With $r(v)=\pi(v)/6$, the normalised gaps
$|\delta|$ over the undirected edges are
\[
  C_6:\quad \underbrace{\tfrac16,\tfrac16,\tfrac16,\tfrac16,\tfrac16}_{\text{edges }(0,1),\dots,(4,5)},\ \underbrace{\tfrac56}_{(5,0)}
  \qquad\Big(\textstyle\sum=\tfrac{10}{6}\Big);
\]
\[
  C_3\sqcup C_3:\quad \underbrace{\tfrac16,\tfrac16}_{(0,1),(1,2)},\ \underbrace{\tfrac26}_{(0,2)},\ \underbrace{\tfrac16,\tfrac16}_{(3,4),(4,5)},\ \underbrace{\tfrac26}_{(3,5)}
  \qquad\Big(\textstyle\sum=\tfrac{8}{6}\Big).
\]

Now fix parameters for a single layer on one channel. Take the node encoder to
output the constant $0$; then $h_v^{(0)}=0$ and the residual, if present,
contributes $0$ (it is the identity on equal widths). Let
$\operatorname{pre}\equiv 1$ (a constant map: zero the output-layer weight and set
its bias to one), so $m_{uv}=1$. Realise the gate $g(\delta)=1+|\delta|$: the
first layer of $\operatorname{MLP}_g$ selects the $|\delta|$ coordinate of
$[\delta,|\delta|,\operatorname{sign}\delta]$, the $\mathrm{ReLU}$ leaves it
unchanged since $|\delta|\ge 0$, and the second layer copies it with weight one.
Let $\operatorname{post}$ be the two-layer MLP computing the identity on the
non-negative achievable values (both layers weight one; the intervening
$\mathrm{ReLU}$ acts as the identity because the values are non-negative). Every
intermediate value is non-negative, so the block activation acts as the identity
too (\cref{remark:bn}). Each node then has
\[
  h_v'=\sum_{u\in N(v)} g(\delta_{uv})\,m_{uv}=\sum_{u\in N(v)}\big(1+|\delta_{uv}|\big),
\]
and, writing $m$ for the number of \emph{undirected} edges (so $E$ contains $2m$
directed edges), the \textsf{add} readout returns, summing over directed edges
(each undirected edge in both orientations, with $|\delta_{uv}|=|\delta_{vu}|$),
\[
  \text{out}(G)=\sum_{v}\sum_{u\in N(v)}\big(1+|\delta_{uv}|\big)
  =2m+2\!\!\sum_{\text{undirected }uv}\!\!|\delta_{uv}| .
\]
Both graphs have $m=6$ undirected edges, so
\[
  \text{out}(C_6)=12+\tfrac{20}{6}=\tfrac{92}{6},
  \qquad
  \text{out}(C_3\sqcup C_3)=12+\tfrac{16}{6}=\tfrac{88}{6},
\]
which differ. LDNA with this canonical rank therefore distinguishes the
1-WL-indistinguishable pair $(C_6,\,C_3\sqcup C_3)$, and by \cref{def:expr} it is
strictly more expressive than 1-WL.
\end{proof}

The theorem asserts the \emph{existence} of a canonical labelling under which
LDNA separates the pair; fixing representatives is legitimate precisely because
a canonizer is a choice of one representative per isomorphism class.

\subsection{Closing remarks}
\label{sec:expr-closing}

\begin{remark}[The \textsf{gru} readout]
\label{remark:gru}
The \textsf{gru} readout runs a recurrent network over the sorted node sequence
and is order-sensitive. Under the feature sort, the order within a block of tied
(identical-row) nodes is inherited from the arbitrary input labelling, so on
graphs with ties the resulting graph function is not permutation-invariant and
is excluded from the statements above. On tie-free graphs (all rows distinct)
there is no expressiveness to be gained either: with an injective initial
colouring, 1-WL is complete after one round. Indeed, the round-$1$ colour of a
node is $(x_v,\{\!\{x_u:u\in N(v)\}\!\})$; since the features are distinct
identifiers, matching two graphs' nodes by feature is a bijection, and equal
round-$1$ colour multisets force equal neighbour-feature sets at every node,
i.e.\ equal labelled adjacency, hence isomorphism. So any two non-isomorphic
tie-free graphs are already separated by 1-WL.
\end{remark}

\begin{remark}[A trilemma]
\label{remark:trilemma}
Informally, and as a synthesis of the three results above rather than a separate
theorem: of the three properties --- deterministic exact invariance,
individualization beyond 1-WL, and cheap worst-case preprocessing --- a rank
scheme can secure at most two. Arbitrary or random tie-breaking (random node
individualization)
attains the extra expressiveness but is invariant only in distribution
\citep{abboud2021rni}; invariant tie-breaks keep exact invariance but fail on
vertex-transitive pairs (\cref{prop:invariant-obstruction}); canonical
individualization keeps both invariance and expressiveness at the cost of a
canonization step (\cref{def:canonical-rank}).
\end{remark}

\begin{remark}[Capacity, not learnability]
\label{remark:learnability}
All results above concern representational capacity --- the existence of
parameters realising a given behaviour. They say nothing about whether such
parameters are found by optimization, nor about generalization.
\end{remark}
